\documentclass[10pt,a4paper]{amsart}
\usepackage[margin=19mm]{geometry}
\usepackage{amsmath,amssymb,mathtools}
\usepackage[scr=rsfs,cal=euler]{mathalfa}
\usepackage{xfrac}
\usepackage[unicode,hidelinks,bookmarks=false]{hyperref}
\newcommand{\ie}{i.e.\ }
\newcommand{\cf}{cf.\ }
\newcommand{\resp}{respectively\ }
\newcommand{\Subsection}{\subsection}
\providecommand{\nobreakdash}{\nobreak\hbox{-}}
\setlength{\emergencystretch}{2em}
\multlinegap0pt
\usepackage{amssymb}
\newtheorem{theorem}{Theorem}
\newtheorem{lemma}{Lemma}
\title{On multiply induced characters}
\author{Masanari Kida, Kyohei Matsukawa, and Hikari Watanabe}
\date{Source: arXiv:2609.04764v1 (CC BY 4.0)}
\begin{document}
\maketitle

\noindent\emph{Source and licence.} On multiply induced characters. Version arXiv:2609.04764v1 (CC BY 4.0), distributed by arXiv under the Creative Commons Attribution 4.0 International licence, http://creativecommons.org/licenses/by/4.0/, as recorded in the arXiv licence metadata for this exact version and shown on the abstract page https://arxiv.org/abs/2609.04764v1. Full licence notice: \texttt{licenses/arxiv-2609.04764-CC-BY-4.0.txt}.\par\smallskip

\noindent\textit{Editorial scope.} Counted unit: the article's theorem prop:4.5 (source0.tex lines 657--662). The article's complete original proof of it is delivered with it (source0.tex lines 675--691, one \texttt{proof} environment of 17 lines, which the article places away from the statement). No mathematics is written, repaired or re-derived: the statement and the proof are the article's own lines, in the article's own notation, and no retained line is substituted or re-flowed.  Retained uncounted as the source-local context the proof needs: lines 666--668.  Nothing was omitted from any retained span, so the package carries no omission record.  No retained line contains a citation, so the wrapper carries no bibliography and no bibliography entry.  No retained line contains a figure environment, an image-inclusion command or drawing code, so no image is delivered and the package declares no asset dependency.  Editorial formatting only, in addition: an AMS \texttt{amsart} preamble that restates the article's own declarations and macros used by the retained lines, namely the environments \texttt{theorem}, \texttt{lemma} on separate counters, together with the standard packages those lines need.  The retained environments print in this document as Theorem 1, Lemma 1 in order of appearance, whereas the article prints the same items with the numbering of its own full document.  The complete native source of the article as distributed in the arXiv e-print for this exact version is bundled unchanged as \texttt{sources/209389/source0.tex}.
\begin{lemma} \label{lem:4.5}
If $(G,N,M)$ is a Camina triple, then $Z (G) \le N$ and $M \le G'$ hold.
\end{lemma}

\begin{theorem}  \label{prop:4.5}
Let $G$ and $H$ are isoclinic groups by the isoclinism $(\eta, \xi )$.
Suppose that $(G,N_1,M_2)$ is a Camina triple of $G$. If we define subgroups $N_2$ and $M_2$ of $H$ by
$\eta (N_1/Z (G)) = N_2 /Z (H)$ and $ \xi (M_1) = M_2$, then $(H,N_2, M_2)$ is a Camina triple 
of $H$.
\end{theorem}

\begin{proof}[Proof of Theorem \ref{prop:4.5}]
    By Lemma \ref{lem:4.5}, $N_2$ is a normal subgroup of $H$ containing $Z (H)$. Also 
    $M_2$ is a normal subgroup of $H$ contained in $H'$ since,
     if $\eta (gZ (G)) = h Z(H) \ (g \in G, h \in H)$, then it is easy to see that
    $\xi (m^g) = \xi (m)^h $ holds for $m \in G'$.

    It remains to prove that, every $h \in H-N_2$ is conjugate to all of $hM_2$.
    We take $g \in G$ satisfying $\eta (g Z(G)) = h Z(G)$. Since $h \not\in N_2$,
    we see that $g \not\in N_1$. As $(G,N_1,M_1)$ is a Camina triple, there exists $x \in G$
    such that $g m_1  = g^x $ for all $m_1 \in M_1$.
This yields $g m_1 Z(G) = g^x Z(G)$. Sending the both sides by $\eta$, we obtain 
$h \eta (m_1) Z (H) = \eta (g)^x Z (H)$. It follows from the definition of 
isoclinism that $\eta (m_1) = \xi (m_1) $ for $m_1 \in G'$. Thus we have
$h \xi (m_1) = \eta (g)^x z = (\eta (g) z)^x$  for some $z \in Z(H)$.
Since $z \in N_2$ and $\eta (g) \not\in N_2$, we see $\eta (g) z \not\in N_2$.
This is what we want to prove.
\end{proof}

\end{document}
